With orthonormal coframe $\theta^0 = d\tau$, $\theta^i = a_i dx^i$, the first structure equation $d\theta^a = -\omega^a{}_b \wedge \theta^b$ determines the connection uniquely (torsion-free, metric). Computing: $d\theta^1 = (\partial_\tau a_1)\,d\tau \wedge dt = \sigma_1\,\theta^0 \wedge \theta^1$, and likewise for $\theta^2$; while $d\theta^3 = \partial_\tau(a_3)\,d\tau\wedge d\varphi\,+\,\partial_\theta(r\sin\theta)\,d\theta \wedge d\varphi = \sigma_3\,\theta^0\wedge\theta^3 + \frac{\cot\theta}{r}\,\theta^2 \wedge \theta^3$. Matching term by term gives $\omega^i{}_0 = \sigma_i\,\theta^i$ (no sum) and $\omega^3{}_2 = \frac{\cot\theta}{r}\,\theta^3$, with all other independent components zero; uniqueness of the Levi-Civita connection closes the argument. Evaluating on the unit vectors and assembling into bivectors gives the display. The last sentence is the standard statement that the induced connection on the constant-$\tau$ spheres is that of $S^2$ of radius $r$, whose coefficients in an orthonormal frame scale as $1/r$; the $\cot\theta$ divergence at the poles is an artifact of the spherical frame, removable by a change of patch.
